\section{Error reduction of close subspaces}

Arad, Landau, Vazirani and Vidick~\cite[Lemma~l:er]{Arad2016RigorousRG}
improve a viable set by applying an approximate ground state projection.
The operator step of their argument holds for any operator that preserves
the target subspace and its orthogonal complement, does not shrink vectors
of the target, and contracts its complement.

\begin{theorem}[Monotonicity of closeness]
    \label{thm:subspace_close_mono}
    \lean{Submodule.IsCloseTo.mono}
    \leanok
    \uses{thm:subspace_close_iff}
    If $U$ is $\delta$-close to $T$ and $U\subseteq U'$, then $U'$ is
    $\delta$-close to $T$.
\end{theorem}

\begin{proof}
    \leanok
    Since $P_U=P_UP_{U'}$, we have
    $\lVert P_Ut\rVert\le\lVert P_{U'}t\rVert$ for every $t$.
\end{proof}

\begin{theorem}[Error reduction by a contraction]
    \label{thm:subspace_close_contract}
    \lean{Submodule.IsCloseTo.map_of_contract}
    \leanok
    \uses{thm:subspace_close_iff, thm:subspace_close_preimage}
    Let $K$ map $T$ and $T^\perp$ into themselves, with
    $\lVert Kt\rVert\ge\lVert t\rVert$ for $t\in T$ and
    $\lVert Kq\rVert^2\le\Delta\lVert q\rVert^2$ for $q\in T^\perp$, where
    $\Delta\ge0$. If $U$ is $\delta$-close to $T$ with $\delta<1$, then $KU$
    is $\Delta/(1-\delta)^2$-close to $T$. This is the operator step in the
    proof of Lemma~l:er of~\cite[lines 606--622]{Arad2016RigorousRG}: for
    $K=\sum_iL_i\otimes A_i\otimes R_i$ and
    $U=\mathcal H_L\otimes S\otimes\mathcal H_R$, the image $KU$ lies in
    $\mathcal H_L\otimes V\otimes\mathcal H_R$ with
    $V=\operatorname{span}\{A_iS\}$. The tensor-product statement of
    Lemma~l:er, with $T$ a spectral subspace of a Hamiltonian, is not
    formalized here.
\end{theorem}

\begin{proof}
    \leanok
    For $t\in T$, the restriction of $K$ to $T$ is injective, hence
    surjective, so $t=Kr$ with $r\in T$ and $\lVert r\rVert\le\lVert t\rVert$.
    Take $u\in U$ with $P_Tu=r$ and
    $\lVert u\rVert\le\lVert r\rVert/(1-\delta)$, and put $q=u-r\in T^\perp$.
    Then $Ku=t+Kq$ with orthogonal summands, and
    $y=\lVert Kq\rVert^2\le\Delta\lVert u\rVert^2
    \le\Delta\lVert t\rVert^2/(1-\delta)^2$. Since $Ku\in KU$,
    $\lVert t\rVert^2=\langle Ku,P_{KU}t\rangle
    \le\lVert Ku\rVert\,\lVert P_{KU}t\rVert$, so with
    $x=\lVert t\rVert^2$ we get
    $\lVert P_{KU}t\rVert^2\ge x^2/(x+y)\ge x-y$.
\end{proof}
